\section{Proofs of the retained decision-theoretic results}
\label{app:b-proofs}

\Cref{as:baseline,as:unit,as:frozen,as:continuity} are in force as stated in
\cref{sec:b-theory}; each proof names the ones it uses.  Throughout,
expectations conditional on $J=j$ are taken over the population of eligible
checkpoints in that stratum, and $\Gset(H_i)$ is identified with
$\Gset(M_i)=\{g:A_{ig}=1\}$.  All regime sets are finite, so pointwise
maximisers admit measurable selections (take the least index among maximisers).

The ordinal-sufficiency theorem is the known advantage-threshold result of the
budgeted-triage literature restated ordinally (Okati et al., 2021;
Fern\'andez-Lor\'ia and Provost, 2022; Kamran et al., 2024); the proof is
retained for completeness and the paper does not claim it as new.  The
certification proposition is stated at the scope actually planned for the
first study; the historical $N=291$ family is not inherited.

\subsection{Proof of \texorpdfstring{\cref{thm:ordinal}}{the ordinal-sufficiency theorem}}

Throughout this proof $W=(Z,M)$ ranges over the common-availability domain and
both scores are read as functions of $w$; the argument uses only that $a_v$ and
$r$ are real-valued measurable scores on a common probability space, so the
special case $r(w)=r(z)$ of the registered risk score is covered without change.

Let $\tau$ satisfy $\Prb(a_v(W)\ge\tau)=q$; such a $\tau$ exists and
$\Prb(a_v(W)=\tau)=0$ by \cref{as:continuity}.  Write
$S^*=S_q(a_v)=\{w:a_v(w)\ge\tau\}$ and let $S$ be any measurable set with
$\Prb(W\in S)=q$; continuity of the allocating score guarantees that
$S=S_q(r)$ is such a set.

\paragraph{Optimality of the upper level set.}
Decomposing over the symmetric difference,
\begin{equation}
 \int_{S^*}a_v\,d\Prb-\int_{S}a_v\,d\Prb
 =\int_{S^*\setminus S}a_v\,d\Prb-\int_{S\setminus S^*}a_v\,d\Prb .
 \label{eq:symdiff}
\end{equation}
On $S^*\setminus S$ we have $a_v\ge\tau$, and on $S\setminus S^*$ we have
$a_v<\tau$.  Moreover $\Prb(S^*\setminus S)=\Prb(S^*)-\Prb(S^*\cap S)
=q-\Prb(S^*\cap S)=\Prb(S)-\Prb(S\cap S^*)=\Prb(S\setminus S^*)$.  Hence
\begin{equation}
 \int_{S^*}a_v\,d\Prb-\int_{S}a_v\,d\Prb
 \ \ge\ \tau\,\Prb(S^*\setminus S)-\tau\,\Prb(S\setminus S^*)=0 .
 \label{eq:np}
\end{equation}
Dividing by $q$ and applying this with $S=S_q(r)$ gives
$\mathcal R_{\mathrm{rank}}(q)\ge0$.

\paragraph{The equality case.}
Rewriting \eqref{eq:symdiff} and \eqref{eq:np},
\begin{equation}
 \int_{S^*}a_v\,d\Prb-\int_{S}a_v\,d\Prb
 =\int_{S^*\setminus S}(a_v-\tau)\,d\Prb+\int_{S\setminus S^*}(\tau-a_v)\,d\Prb,
\end{equation}
both integrands being non-negative.  So $\mathcal R_{\mathrm{rank}}(q)=0$ forces
$a_v=\tau$ almost everywhere on $S^*\setminus S$ and on $S\setminus S^*$.
Continuity of the law of $a_v(W)$ gives $\Prb(a_v(W)=\tau)=0$, hence
$\Prb(S^*\setminus S)=\Prb(S\setminus S^*)=0$: the two selected sets coincide up
to a null set.

If this holds for every $q\in(0,1)$ then $S_q(r)\doteq S_q(a_v)$ for every $q$,
so $r$ and $a_v$ have the same family of upper level sets up to null sets, which
is exactly the statement that they induce the same ordering almost everywhere.
The converse is immediate: identical orderings give identical upper level sets
at every capacity, hence equality at every $q$. \qed

\begin{remark}[Association versus ordinal sufficiency]
\label{rem:association}
Let $s$ be any monitor score and $\Delta_v$ the verification advantage on the
common domain, both with continuous laws.  (i)~If $s$ and $\Delta_v$ induce the
same ordering almost everywhere --- the null of \cref{thm:ordinal} --- then
$\operatorname{Cov}(s,\Delta_v)\ge0$, since comonotone random variables have
non-negative covariance.  (ii)~The converse fails, and fails exactly where the
budget is spent: for any $q\in(0,1)$ there is a joint law with
$\operatorname{Cov}(s,\Delta_v)>0$ and $\mathcal R_{\mathrm{rank}}(q)>0$.  Give
the top-$q$ mass of $s$ an advantage concentrated near zero and set
$\Delta_v=s$ on the remaining $1-q$; the covariance stays positive, while the
top-$q$ set under $s$ is not the top-$q$ set under $\Delta_v$, so the advantage
mass lost on the selected fraction, $q\,\mathcal R_{\mathrm{rank}}(q)$, is
positive.  (iii)~Even the perturbation $\Delta_v=s+\varepsilon N$ with $N$
independent has $\operatorname{Cov}(s,\Delta_v)=\operatorname{Var}(s)>0$ for
every $\varepsilon>0$ while $\mathcal R_{\mathrm{rank}}(q)>0$ at every $q$, so a
merely noisy but monotone monitor is not ordinally sufficient either.  Hence a
return-association gate of the kind used to weight privileged credit
\citep[Prop.~1]{adrs2026} is implied by, but does not imply, the ordering E3
tests, and it says nothing about the top-$q$ set that a fixed verification count
actually spends.
\end{remark}

With $2\times10^5$ seeded draws and $s\sim\mathrm{Unif}(0,1)$: construction~(ii)
at $q=0.05$ has $\operatorname{Cov}\approx0.060$, and the advantage mass lost on
the selected fraction is $q\,\mathcal R_{\mathrm{rank}}(q)\approx0.046$ at
$q=0.05$ and $\approx0.024$ at $q=0.5$; construction~(iii) at
$\varepsilon=0.05$ has
$\operatorname{Cov}\approx0.083=\operatorname{Var}(s)$ and
$q\,\mathcal R_{\mathrm{rank}}(q)\approx1.2\times10^{-3}$ at each of the four
registered capacities of \eqref{eq:main-capacities}; the comonotone control
gives $\mathcal R_{\mathrm{rank}}\equiv0$.  These values are reproduced in
\texttt{verify\_theory.py}.

\subsection{Verification of the construction in \texorpdfstring{\eqref{eq:construction}}{the construction equation}}

Let $Z\sim\mathrm{Unif}[0,1]$ with the mask degenerate, so that $W$ is identified
with $Z$, and set $r(z)=z$ and $a_v(z)=\tfrac12-z$; both laws are
continuous, so \cref{thm:ordinal} applies.  Selecting the upper $q$-fraction by
$r$ gives $S_q(r)=[1-q,1]$, on which $\E[Z]=1-\tfrac q2$, so
\begin{equation}
 \E\big[a_v(Z)\mid Z\in S_q(r)\big]=\tfrac12-\big(1-\tfrac q2\big)=\tfrac q2-\tfrac12 .
\end{equation}
Selecting by $a_v$ instead favours small $z$, giving $S_q(a_v)=[0,q]$, on which
$\E[Z]=\tfrac q2$, so
$\E[a_v(Z)\mid Z\in S_q(a_v)]=\tfrac12-\tfrac q2$.  Subtracting,
\begin{equation}
 \mathcal R_{\mathrm{rank}}(q)=\Big(\tfrac12-\tfrac q2\Big)-\Big(\tfrac q2-\tfrac12\Big)=1-q,
\end{equation}
strictly positive for every $q\in(0,1)$ and approaching the maximal attainable
regret as the verification capacity shrinks.  Replacing $a_v$ by $z-\tfrac12$
leaves the law of $r(Z)$ unchanged and yields
$\mathcal R_{\mathrm{rank}}\equiv0$.  The two populations are therefore
indistinguishable by any functional of the risk-score distribution alone, and
differ maximally in whether that score is usable for allocation.

\subsection{Proof of \texorpdfstring{\cref{thm:refinement}}{the refinement theorem}}

\paragraph{Measurability of the feasible set.}
By definition $M_i=(A_{ig}:g\in\Gset)$ with
$A_{ig}=\mathbf 1\{g\in\Gset(H_i)\}$, so $\Gset(H_i)$ is a deterministic
function of $M_i$.  Hence $\Gset(H_i)$ is $\sigma(Z_i,M_i)$-measurable, and the
maximisations below range over a set both policy classes can identify.

\paragraph{The two suprema.}
Let $d\in\Dcal_H$.  Since $d(H)$ is $H$-measurable and takes values in
$\Gset(H)$, the tower property gives
\begin{equation}
 V_j(d)=\E\big[Y^{d(H)}\mid J=j\big]
       =\E\big[\mu^{(j)}_{d(H)}(H)\mid J=j\big]
       \le\E\Big[\max_{g\in\Gset(H)}\mu^{(j)}_g(H)\;\Big|\;J=j\Big].
\end{equation}
The bound is attained by the $H$-measurable selector
$d(H)\in\argmax_{g\in\Gset(H)}\mu^{(j)}_g(H)$, which establishes the expression
for $V^{*H}_j$.  For $d\in\Dcal_Z$ the identical argument, conditioning on
$(Z,M)$ rather than $H$, gives
$V_j(d)=\E[\mu^{(j)}_{d(Z,M)}(Z,M)\mid J=j]$ and the expression for $V^{*Z}_j$;
here the attaining selector is $(Z,M)$-measurable precisely because $\Gset(H)$
is, by the previous paragraph.

\paragraph{Non-negativity of the refinement loss.}
Since $\sigma(Z,M)\subseteq\sigma(H)$, the tower property gives
$\mu^{(j)}_g(Z,M)=\E[\mu^{(j)}_g(H)\mid Z,M,J=j]$ for every $g$ available at the
checkpoint.  For each such $g$,
\begin{equation}
 \mu^{(j)}_g(H)\le\max_{g'\in\Gset(H)}\mu^{(j)}_{g'}(H)
 \quad\text{pointwise,}
\end{equation}
so taking $\E[\,\cdot\mid Z,M,J=j]$ on both sides,
\begin{equation}
 \mu^{(j)}_g(Z,M)\le\E\Big[\max_{g'\in\Gset(H)}\mu^{(j)}_{g'}(H)\;\Big|\;Z,M,J=j\Big].
\end{equation}
The right-hand side does not depend on $g$, so the inequality survives taking
the maximum over $g\in\Gset(H)$ on the left --- this is the conditional Jensen
inequality for the convex map $\max$.  Taking $\E[\,\cdot\mid J=j]$ and applying
the tower property once more yields $V^{*Z}_j\le V^{*H}_j$, i.e.
$\mathrm{Ref}_j\ge0$.

\paragraph{Sufficiency.}
If the stratum-wise form of \eqref{eq:conditional-mean-sufficiency} holds, then
$\mu^{(j)}_g(H)=\mu^{(j)}_g(Z,M)$ almost surely for every available $g$.  The
two maxima therefore coincide almost surely, and $\mathrm{Ref}_j=0$.

\paragraph{The estimable lower bound.}
If $\sigma(Z)\subseteq\sigma(Z')\subseteq\sigma(H)$ then every $(Z,M)$-measurable
policy is $(Z',M)$-measurable and every $(Z',M)$-measurable policy is
$H$-measurable, so $\Dcal_Z\subseteq\Dcal_{Z'}\subseteq\Dcal_H$ and the suprema
are ordered: $V^{*Z}_j\le V^{*Z'}_j\le V^{*H}_j$.  Subtracting $V^{*Z}_j$ gives
\eqref{eq:refinement-bound}. \qed

\medskip
Only the last display is estimable: $V^{*Z'}_j-V^{*Z}_j$ compares two policy
classes whose members can both be evaluated on held-out branches, whereas
$V^{*H}_j$ requires conditional means at checkpoint resolution.  This asymmetry
is the reason \eqref{eq:conditional-mean-sufficiency} admits falsification but
not confirmation.

\begin{remark}[Where the conditional expectations may not be dropped]
\label{rem:aggregation}
Both suprema in \cref{thm:refinement} are expectations of a maximum of
\emph{conditional means}, and the Bellman reading of \cref{sec:bellman} does not
license replacing either by an expectation of a maximum of the state-level
action value.  Writing $\mu^{(j)}_g(H)=\E[Q^\rho_\tau(X,g)\mid H,J=j]$, the same
conditional Jensen step used above gives
\begin{equation}
 V^{*Z}_j\;\le\;V^{*H}_j
 =\E\Big[\max_{g\in\Gset(H)}\E\big[Q^\rho_\tau(X,g)\mid H,J=j\big]\Big]
 \;\le\;\E\Big[\max_{g\in\Gset(H)}Q^\rho_\tau(X,g)\;\Big|\;J=j\Big],
 \label{eq:aggregation-chain}
\end{equation}
and the last inequality is generally strict: the right-hand side is the value of
a policy that could see the residual randomness in $X$ beyond $H$, which no
admissible policy can.  Only the two left-hand quantities are policy values in
$\Dcal_H$ and $\Dcal_Z$, and only their difference is what E1 estimates.  The
same point applies one level down: the actionable value \eqref{eq:actionable} is
$\max_g$ of conditional means and is \emph{not} $\E[\max_g A^\rho\mid Z,M,J]$.
A dictionary row that dropped the inner conditional expectation would silently
substitute an unattainable oracle for the estimand.
\end{remark}

\subsection{Proof of \texorpdfstring{\cref{prop:certification}}{the certification proposition}}

There are $|\Dcal|$ policies and, for each, $K+3$ bounds: one for the value used
in selection, one for the no-fault contrast of \eqref{eq:noninferiority}, one
for the invalid-output rate of \eqref{eq:invalid-constraint}, and one for each of
the $K$ cost components of \eqref{eq:cost-constraint}.  Together with the gate
family of \eqref{eq:gate-family}, the combined family has $N=|\Dcal|(K+3)+F_{\mathrm{gate}}$
members.  Allocating $\delta/N$ to each and applying a union bound, the
probability that any bound in the combined family fails is at most $\delta$;
in particular every certification bound is valid, and every gate bound is valid
(which is exactly the uniform-validity hypothesis of \cref{thm:safe}).

On the complement $\mathcal E'$, every certification bound is valid.  A policy
$d$ that is certified --- its bounds place it inside every constraint ---
therefore has its true quantities inside every constraint on $\mathcal E'$, so
$d\in\mathcal C$; holding for all $d$ at once gives
$\widehat{\mathcal C}\subseteq\mathcal C$, and any policy selected from
$\widehat{\mathcal C}$ by any rule --- including one that depends on which
constraints certified --- satisfies every constraint.  It is this last clause
that requires simultaneity over $\Dcal$: without it, $\widehat{\mathcal C}$ is a
random set whose selection event is not controlled by per-policy coverage.

For the half-width, let member $k$ be the mean of $n_k$ independent
contributions $X_1,\dots,X_{n_k}$, each in an interval of length $R_k$.
Hoeffding's inequality gives
$\Prb(|\widehat\theta_k-\theta_k|\ge\varepsilon)\le
2\exp(-2n_k\varepsilon^2/R_k^2)$.  Setting the right-hand side equal to
$\delta/N$ and solving gives \eqref{eq:halfwidth}; certification of a
constraint with true slack $s_k$ requires $\varepsilon_k(n_k)\le s_k$, which
rearranges to \eqref{eq:sample-ceiling}.  For the clustered form, let the
member be $\sum_c w_c m_c$ over independent clusters $c$ with $\sum_c w_c=1$
and each $m_c$ in an interval of length $R_k$.  The summands $w_cm_c$ are
independent with ranges $w_cR_k$, so Hoeffding gives
$\Prb(|\sum_c w_c(m_c-\E m_c)|\ge\varepsilon)\le
2\exp\!\big(-2\varepsilon^2/(R_k^2\sum_c w_c^2)\big)$, and solving at
$\delta/N$ yields $\varepsilon_k=R_k\sqrt{\log(2N/\delta)\sum_cw_c^2/2}$.  With
$w_c=1/n_k$ this is \eqref{eq:halfwidth}; with a single cluster
$\sum_cw_c^2=1$ and the width is at least $R_k\sqrt{\log(2N/\delta)/2}$, which
exceeds $R_k$ for every $N\ge 2$ and $\delta\le 1/2$, so one cluster never
certifies whatever its size. \qed

\medskip
At $\delta=0.05$, $K=8$, $|\Dcal|=9$, $F_{\mathrm{gate}}\le 192$ and $s=0.05$, the
combined family has $N=291$ members and the argument of the logarithm is
$2\cdot291/0.05=1.164\times10^{4}$, so with $R_k=1$ and $n_k=n$ the historical
unit-range, equal-count, independent-unit illustration gives
$n\ge 9.362/(2\times0.0025)=1.87\times10^{3}$.  That figure is \emph{superseded}
as a planning number (author instruction A9, 2026-09-23): the body states no
bound-derived count, and the corrected stratum-wise recomputation under the
frozen selection split is recorded in the analysis plan.  The derivation is kept
here only to show how the old illustration behaved.  For the no-fault member,
$R_k=2$ and the same slack needs $R_k^2=4$ times the unit-range count of
independent units; for a gate member on one cell and stratum, $n_k$ is that
cell's own count.  The bound
is conservative in variance, so a variance-adaptive construction with the
pilot's observed variance will be smaller, but it is not conservative in range
or in unit: a unit-range figure quoted for a range-two member, or a root count
quoted for a template count, understates the requirement.

\subsection{Proof of \texorpdfstring{\cref{thm:safe}}{the safe-branching theorem}}

Work on the event $\mathcal E$, of probability at least $1-\delta$, on which
every bound is valid.

By the tower property, for any $(Z,M)$-measurable policy $d$,
\begin{equation}
 V_j(d)=\E\big[\mu^{(j)}_{d(Z,M)}(Z,M)\mid J=j\big],
\end{equation}
so that, using \eqref{eq:pointwise-advantage},
\begin{equation}
 V_j(d^\dagger)-V_j(d^\star)
 =\E\big[\Delta^{(j)}_{d^\dagger(Z,M)}(Z,M)\mid J=j\big].
 \label{eq:one-step-pd}
\end{equation}
This is the one-step performance-difference identity; it is exact and consists
of a single term because the decision is taken once.

Partition the eligible support by the deviation set
\begin{equation}
 \mathcal N=\Big\{(z,m):\max_{g:A_g=1}\ \min_{j'\in\{0,1\}}
 \mathrm{LCB}_\delta\big[\Delta^{(j')}_g(c(z,m))\big]>0\Big\}.
\end{equation}
All bounds are the cell-level bounds of \eqref{eq:cell-advantage}; with
cell-measurable policies the integrand $\Delta^{(j)}_{d(Z,M)}(Z,M)$ equals the
cell-level advantage of the regime selected at $c(Z,M)$, so cell-level validity
applies pointwise in the integrand.

Off $\mathcal N$, rule \eqref{eq:gated-rule} sets $d^\dagger(z,m)=d^\star(z,m)$,
and by \eqref{eq:pointwise-advantage},
$\Delta^{(j)}_{d^\star(z,m)}(z,m)=\mu^{(j)}_{d^\star(z,m)}(z,m)-
\mu^{(j)}_{d^\star(z,m)}(z,m)=0$ identically.  The integrand contributes
exactly zero; this is the consequence of defining the advantage against the
comparator rather than against $g_0$.

On $\mathcal N$, let $g^\dagger(z,m)$ be the selected regime.  It attains a
positive value of $\min_{j'}\mathrm{LCB}_\delta[\Delta^{(j')}_{g^\dagger}]$, so
for \emph{each} $j\in\{0,1\}$,
\begin{equation}
 \Delta^{(j)}_{g^\dagger}(c(z,m))\ \ge\ \mathrm{LCB}_\delta\big[\Delta^{(j)}_{g^\dagger}(c(z,m))\big]\ >\ 0
 \qquad\text{on }\mathcal E,
\end{equation}
the first inequality by validity and the second by the gate.  Taking the
minimum over $j$ before gating is what makes this hold in both strata rather
than in one.

Every value of the integrand in \eqref{eq:one-step-pd} is therefore
non-negative, giving $V_j(d^\dagger)\ge V_j(d^\star)$ for $j=0,1$.  Finally, for
any $p\in[0,1]$,
\begin{equation}
 V(d^\dagger;p)-V(d^\star;p)
 =p\big\{V_1(d^\dagger)-V_1(d^\star)\big\}+(1-p)\big\{V_0(d^\dagger)-V_0(d^\star)\big\}\ \ge\ 0,
\end{equation}
since both weights and both brackets are non-negative.  The deviation set
$\mathcal N$ does not depend on $p$, so a single event $\mathcal E$ carries the
conclusion for every $p$ at once. \qed

\begin{remark}[Where each hypothesis is used]
Freezing $d^\star$ before the bounds are constructed keeps the target of
inference fixed: $\Delta^{(j)}_g$ has a population value whether or not the
comparator was adapted, and the defect of reusing calibration data is not that
the quantity stops being a population quantity but that the coverage statement
then refers to a data-dependent estimand, voiding the declared selection and
certification argument.  Uniform validity over the eligible support licenses the
partition: a bound valid only marginally would not survive conditioning on
$d^\dagger$ having chosen to deviate.  Under full-arm execution every feasible
regime is executed at every sampled checkpoint, so within the cells the
sampling law reaches the \emph{coverage} half is a property of the design ---
provided restoration and isolation hold, which is the instrument gate; a cell
with no sampled checkpoints has no bound and the gate falls back there.  The \emph{statistical} half is not: uniform validity of the gate family
of \cref{eq:gate-family} is established on held-out data by a union bound over
the registered finite partition, union-bounded together with the certification
family of \cref{prop:certification}, and the cell count $B$ enters the combined
ceiling \eqref{eq:sample-ceiling}.
\end{remark}
\begin{remark}[What is not guaranteed]
\Cref{thm:safe} bounds the completion probability and nothing else.  It does not
bound invalid-output rates, cost, or the number of checkpoints at which an
individual regime is harmful; those are separate constraints with their own
tolerances, certified under \cref{prop:certification} and measured directly.
\end{remark}

\subsection{Proof of \texorpdfstring{\cref{cor:interval}}{the guarantee-interval corollary}}

$\Gamma(p)=p\Gamma_1+(1-p)\Gamma_0=\Gamma_0+p(\Gamma_1-\Gamma_0)$ is affine in
$p$.  The set on which an affine function of one variable is non-negative is a
closed half-line, all of $\mathbb R$, or empty; intersecting with $[0,1]$ gives a
closed subinterval.  If $\Gamma_0$ and $\Gamma_1$ have strictly opposite signs
then $\Gamma_1\neq\Gamma_0$ and the unique root of $\Gamma$ is
\begin{equation}
 p=\frac{-\Gamma_0}{\Gamma_1-\Gamma_0},
\end{equation}
which lies in $(0,1)$, and is the break-even prevalence of
\cref{app:estimands} with $G$ replaced by $\Gamma$.  Under the gate of
\eqref{eq:gated-rule}, \cref{thm:safe} gives $\Gamma_0\ge0$ and $\Gamma_1\ge0$,
so $\Gamma(p)\ge0$ throughout $[0,1]$. \qed

\subsection{Proof of \texorpdfstring{\cref{cor:proposal}}{the proposal corollary}}

Work on the event $\mathcal E$, of probability at least $1-\delta$, on which
every bound is valid.  By (b) and (c), $d_{\mathrm{prop}}(z,m)\in\Gset(z,m)$ at
every eligible checkpoint, so $\Delta^{(j)}_{d_{\mathrm{prop}}(z,m)}(c(z,m))$ and
its lower confidence bound are defined wherever the checkpoint is eligible.
Partition the eligible support by the proposal-deviation set
\begin{equation}
 \mathcal N_p=\Big\{(z,m):\min_{j'\in\{0,1\}}\mathrm{LCB}_\delta
 \big[\Delta^{(j')}_{d_{\mathrm{prop}}(z,m)}(c(z,m))\big]>0\Big\}.
\end{equation}
Off $\mathcal N_p$, rule \eqref{eq:proposal-rule} returns $d^\star$, so by
\eqref{eq:pointwise-advantage} the integrand in the one-step identity
\eqref{eq:one-step-pd} is zero.  On $\mathcal N_p$, rule \eqref{eq:proposal-rule}
deviates to $d_{\mathrm{prop}}(z,m)$, and for each $j\in\{0,1\}$,
\begin{equation}
 \Delta^{(j)}_{d_{\mathrm{prop}}(z,m)}(c(z,m))\ \ge\
 \mathrm{LCB}_\delta\big[\Delta^{(j)}_{d_{\mathrm{prop}}(z,m)}(c(z,m))\big]\ >\ 0
 \qquad\text{on }\mathcal E,
\end{equation}
the first inequality by validity and the second by the gate.  Every value of the
integrand is therefore non-negative in both strata, giving
$V_j(d^{\dagger}_{\mathrm{prop}})\ge V_j(d^\star)$ for $j=0,1$, and the 
prevalence-weighting step of the proof of \cref{thm:safe} carries the conclusion
to every $p\in[0,1]$ at once, since $\mathcal N_p$ does not depend on $p$. \qed

\medskip
Hypothesis (a) is not used in this argument: on a validity event covering every
available regime, the selection rule by which the proposal was chosen is
irrelevant.  It is a protocol condition --- the validity event is constructed on
held-out data over a finite registered class, so a proposal that consumes
calibration feedback produces candidates the registered allocation does not
count, and an engine emitting a fresh arbitrary patch at deployment time is not
covered even though the proof would go through for a fixed one.

\begin{remark}[Two legitimate constructions of a proposal policy]
Option A --- a \emph{frozen patch library}: the engineer generates patches on
the registered pair-development subset only; validated patches are frozen into a finite regime set
that enters branching and certification like any other $g$, so (b) holds by
construction.

Option B --- \emph{the frozen generator as one stochastic regime}: define
$g_{\mathrm{HR1}}$ as ``invoke the frozen engineer under fixed decoding,
validator and response rules''.  Individual patches may differ; the randomness
is absorbed into $Y^{g_{\mathrm{HR1}}}$, estimated by repeated branches.  This is
clean only with all five: a registered checkpoint-level input contract; fixed
decoding and validator rules; pre-treatment availability; a frozen model and
version; and repeated branch evaluation.

The first condition of Option B is the binding one: a generator such as the
harness engineer of \cref{sec:related} consumes a \emph{batch of failures},
whereas our regime fires at a \emph{single checkpoint}.  Without an adapter that
restates its input contract at checkpoint granularity, only its hook design
transfers --- the recovery routine and the guardrail as one-step regimes --- not
the published generator; building that adapter is research, not integration.
\end{remark}

